四问的共同困难不是求一个无误差交点，而是在观测误差、接收半径未知和定向遮蔽下维持一个必然包含真源的集合，并据此作出有限时间内可执行的移动决策。为避免后续把经验现象误当作保证，先将各问的输入、核心决策和验证方式分开。

问题一属于集合反演：每次示向给出一个角度区间，定位区域是全部示向楔形的交集。应先判断交集是否为空或无界，再求有界凸多边形的直径；清除可行性则由最小包围圆半径判断，不能仅凭直径作结论。

问题二属于鲁棒选址：第二点既要对所有可能源位置保证仍在最小接收半径内，又要形成足够大的交会角。本文先用最坏情形距离得到保证接收候选域，再在该域内选择横向基线较大的对称点，把“收到信号”和“定位稳定”两个要求分开处理。

问题三可拆为外层覆盖搜索和内层单源定位。外层需要一个与源数、接收半径实际值无关的有限覆盖点集；内层只用正示向不断相交可行域，并以最小包围圆给出安全清除点。若自适应定位未在限定次数内收敛，则用有限网格覆盖剩余可行域，从而保留完备性。

问题四的关键变化是无信号具有二义性：它既可能表示距离过远，也可能表示检测点位于定向源背面。因此无信号不能用于删减位置集合。需要重新证明同一覆盖点集对任意辐射半平面仍至少存在一个可接收见证点，发现后再复用问题三的定位与清除过程。

四问的验证分别采用几何反例、解析极值与离散收敛、固定种子合成模拟、官方演练日志。时间评价统一按接口协议分解为移动、测量、频道切换和清除四部分；正式测试数据与演练数据严格分栏记录。
